\Needspace{8\baselineskip}
\subsection{Generic uniqueness}
\label{sec:generic-adjoints}

On every affine line in $C(X)$, the concave function $B_\varphi$ is differentiable outside a countable set, so the gap between its two directional derivatives along the line vanishes on a dense subset. Upper semicontinuity of this gap and separability of $C(X)$ turn the one-dimensional fact into uniqueness of the supporting probability on a dense $G_\delta$ subset of $C(X)$.

\Needspace{12\baselineskip}
\begin{theorem}\label{thm:generic-adjoint-uniqueness}\label{cor:generic-adjoint-approximation}\index{Adjoint measure!generic uniqueness}
The set
\[
 \mathscr U_\varphi
   =\{F\in C(X,\mathbb R):\mathcal A_\varphi(F)
                                  \text{ is a singleton}\}
\]
is a dense $G_\delta$ subset of $C(X)$ with the uniform norm. At every $F\in\mathscr U_\varphi$, the functional $B_\varphi$ is G\^ateaux differentiable, and its derivative is minus the integral against the unique supporting probability. If $F_j\to F$ uniformly and $\nu_j\in\mathcal A_\varphi(F_j)$, then $\nu_j$ converges weakly to that probability.
\end{theorem}
\Needspace{5\baselineskip}
\begin{proof}
For $h\in C(X,\mathbb R)$ define the width
\[
 w_h(F)=\max_{\nu\in\mathcal A_\varphi(F)}\int_Xh\,\mathrm{d}\nu
         -\min_{\nu\in\mathcal A_\varphi(F)}\int_Xh\,\mathrm{d}\nu\geq0.
\]
The extrema exist by weak compactness. We claim that the width is upper semicontinuous in $F$. Indeed, let $F_j\to F$ and choose maximizing and minimizing probabilities at each $F_j$. Along a subsequence realizing the upper limit of the widths, compactness of the space of probability measures gives weak limits of both choices. By \cref{prop:weak-adjoint-stability}, both limits belong to $\mathcal A_\varphi(F)$, so the difference of their integrals against $h$ is at most $w_h(F)$.

Consequently, for every $\epsilon>0$ the set $\{F:w_h(F)<\epsilon\}$ is open. It is also dense. If $h=0$ there is nothing to prove. Otherwise fix $F$ and an open ball about it. The function
\[
 b(t)=B_\varphi(F+th)
\]
is finite and concave on the real line, so it is differentiable outside an at most countable set: its one-sided derivatives are finite and monotone, and the nonempty open intervals between unequal one-sided derivatives at distinct points are disjoint, so each contains a different rational number. Choose a point of differentiability $t$ so close to zero that $F+th$ lies in the given ball. The derivative formulas of \cref{prop:normalization-directional} give $w_h(F+th)=0$, which proves density.

The compact manifold $X$ is metrizable, so $C(X)$ is separable. Choose a countable uniformly dense family $(h_j)$ in $C(X)$. By the Baire category theorem in the complete metric space $C(X)$, the set
\[
 \bigcap_{j,k\geq1}\{F:w_{h_j}(F)<1/k\}
\]
is a dense $G_\delta$. At every point of this intersection, all supporting probabilities have the same integral against every $h_j$. Since they have unit mass, uniform approximation gives equal integrals against every continuous function, and the Riesz representation theorem shows that the probabilities coincide. Conversely, a singleton has zero width in every direction, so the intersection is exactly $\mathscr U_\varphi$. The formula for the derivative follows from \cref{prop:normalization-directional}, and the approximation assertion from \cref{prop:weak-adjoint-stability}.
\end{proof}
